
A projekt célja, hogy egy modern, felhasználóbarát webalkalmazást fejlesszen ki a tömegközlekedés népszerűsítésére és a buszok használatának egyszerűsítésére. A következő konkrét célkitűzéseket határozzuk meg:

\begin{itemize}
    \item \textbf{Digitális Jegyvásárlás:} Olyan funkció kialakítása, amely lehetővé teszi a felhasználók számára, hogy előre jegyet vásároljanak a kívánt járatra, így csökkentve a készpénzes tranzakciók számát.
    
    \item \textbf{Valós Idejű Buszkövetés:} Integrálni egy valós idejű buszkövető rendszert, amely lehetővé teszi a felhasználók számára, hogy nyomon követhessék a buszok aktuális helyzetét és várható érkezési idejét.
    
    \item \textbf{Menetrendek Elérhetősége:} Biztosítani, hogy a buszmenetrendek minden megállóban könnyen elérhetők legyenek, valamint online platformon is megtekinthetők legyenek.
    
    \item \textbf{Interaktív Térkép:} Fejleszteni egy interaktív térképet, amely lehetővé teszi a felhasználók számára, hogy könnyen navigáljanak a város tömegközlekedési hálózatán, és információkat kapjanak a megállókról és járatokról.
    
    \item \textbf{Adminisztrációs Felület:} Létrehozni egy adminisztrációs felületet, amelyet a busztársaságok használhatnak a menetrendek, útvonalak és alkalmazottak információinak kezelésére, ezzel növelve a működés hatékonyságát.
    
    \item \textbf{Felhasználói Élmény Javítása:} A webalkalmazás tervezése során a felhasználói élmény maximalizálása, beleértve a könnyű navigációt és az intuitív felület kialakítását.
\end{itemize}

Ezek a célkitűzések a webalkalmazás fejlesztésének alapját képezik, és hozzájárulnak a modern tömegközlekedési rendszerek hatékonyabbá tételéhez, valamint a felhasználók utazási élményének javításához.
